\documentclass{article}

% Pacotes
\usepackage[utf8]{inputenc}
\usepackage[brazil]{babel}
\usepackage[letterpaper,top=2.5cm,bottom=2.5cm,left=3cm,right=3cm,marginparwidth=1.75cm]{geometry}
\usepackage{amsmath, graphicx, subcaption}
\usepackage[colorlinks=true, allcolors=blue]{hyperref}

\usepackage{listings}
\usepackage{xcolor}
\usepackage{tabularx} 
\usepackage{array} 
\usepackage[siunitx, RPvoltages]{circuitikz}
\usepackage{amssymb}      % símbolos matemáticos extras
\usepackage{extarrows}  % para usar \xleftrightarrow
\usepackage{url}
\usepackage{hyperref} % O pacote hyperref geralmente já inclui funcionalidades de quebra de URL
\usepackage{booktabs}  % Para top/mid/bottomrule
\usepackage{float}     % Para [H] nas tabelas
\usepackage{booktabs} % Para top/mid/bottomrule
\usepackage{float}    % Para [H]
\usepackage{graphicx} % Para incluir figuras
\usepackage{enumitem} % Para customizar listas (a), (b), (c)




% Definições de estilo para Python
\lstdefinestyle{pythonstyle}{
      language=Python,
      basicstyle=\ttfamily\small,
      keywordstyle=\color{blue}\bfseries,
      commentstyle=\color{gray}\itshape,
      stringstyle=\color{orange},
      numbers=left,
      numberstyle=\tiny,
      stepnumber=1,
      numbersep=10pt,
      backgroundcolor=\color{white},
      frame=single,
      breaklines=true,
      tabsize=4,
      showstringspaces=false,
}


\begin{document}

% Capa personalizada
\begin{titlepage}
      \centering
      \includegraphics[width=4cm]{images/fig_logo_cefetmg.pdf}\par
      \vspace{0.5cm}
      
      {\LARGE Centro Federal de Educação Tecnológica de Minas Gerais \par}
      \vspace{0.2cm}
      
      {\Large Engenharia de Computação \\ Computação Gráfica \par}
      \vspace{4cm}
      
      {\Huge \textbf{Relatório de Desenvolvimento} \par}
      \vspace{0.5cm}
      
      {\LARGE Simulação de um Braço Robótico Articulado \par}
      \vspace{3cm}
      
      {\LARGE João Paulo da Cunha Faria \par}
      {\LARGE Maria Eduarda Teixeira Souza \par}
      \vfill
      
      {\large Divinópolis - MG \par}
      {\large Agosto de 2026 \par}
\end{titlepage}


\section{Introdução}

A modelagem de manipuladores robóticos é um passo fundamental no estudo da robótica, permitindo prever o posicionamento espacial de suas ferramentas. Ao transpor a complexidade da matemática pura para a simplicidade das ferramentas de processamento gráfico, observa-se que as bibliotecas fornecem uma abstração robusta sobre os algoritmos matemáticos.

Para um sistema de braço robótico articulado operando em um plano bidimensional, a posição de cada junta pode ser determinada por meio das equações de cinemática direta, utilizando trigonometria ($\cos(\theta), \sin(\theta)$).

\section{Modelagem Matemática e Primitivas Gráficas}

O projeto computacional exige o cálculo do posicionamento das articulações e sua posterior renderização na tela. Conforme abordado nas aulas da disciplina, há uma divisão clássica entre o Plano Cartesiano e o Plano de Tela (\textit{Screen Space}).

Devido à herança da arquitetura dos monitores CRT (Tubo de Raios Catódicos), a origem de coordenadas $(0,0)$ de matrizes gráficas em linguagens como Python, C e Java situa-se sempre no canto superior esquerdo, e o eixo $Y$ cresce para baixo. Dessa forma, para que a renderização do robô aponte para "cima" na tela, é estritamente necessário subtrair o resultado do cálculo trigonométrico do eixo $Y$ da base, caso contrário, o robô seria desenhado de ponta-cabeça.

Assim, a posição final $(x_f, y_f)$ de uma junta conectada a um elo de comprimento $L$ sob um ângulo $\theta$ em relação à sua origem $(x_i, y_i)$ é adaptada ao \textit{Screen Space} da seguinte forma:

\begin{equation}\label{eq:cinematica_x}
x_f = x_i + L \cos(\theta)
\end{equation}

\begin{equation}\label{eq:cinematica_y}
y_f = y_i - L \sin(\theta)
\end{equation}

\subsection{Posicionamento com Cinemática Direta}

O braço é composto por uma base, um ombro, um cotovelo e uma ferramenta na extremidade (garra). A partir do ponto fixo da base, o cotovelo é posicionado utilizando o comprimento do braço e o ângulo do ombro. Subsequentemente, a mão é posicionada. Um ponto crucial da modelagem é a dependência do ângulo total acumulado: o ângulo do segundo elo (antebraço) deve ser calculado obrigatoriamente utilizando a soma dos ângulos da junta anterior e da junta atual (\texttt{ang\_ombro\_rad + ang\_cotovelo\_rad}).

\subsection{Melhorias Implementadas Além do Escopo Base}

O desafio do Braço Robótico proposto exigia o cálculo das coordenadas de interseção e o uso de funções primitivas básicas do OpenCV (\texttt{cv2.line}, \texttt{cv2.circle}, \texttt{cv2.rectangle}) para projetar a peça na tela. No entanto, visando aprimorar a estética e aproximar a simulação gráfica de um modelo mecânico real, este trabalho implementou melhorias adicionais estilizadas:

\begin{itemize}
    \item \textbf{Engrenagens Paramétricas:} Ao invés de representar as juntas apenas com círculos preenchidos simples, foi desenvolvida a função \texttt{desenhar\_engrenagem}. Esta função utiliza conceitos de divisão angular trigonométrica para iterar e alocar $N$ pequenos "dentes" de engrenagem ao redor da junta principal, conferindo um aspecto de atuador rotativo industrial.
    \item \textbf{Garra Direcional (Atuador de Extremidade):} A extremidade do robô recebeu uma garra composta por um círculo (base) e dois segmentos de reta (dedos). Para que a garra sempre aponte na mesma direção do antebraço, o algoritmo usa \texttt{math.atan2(dy, dx)} — que converte o vetor cotovelo-mão diretamente em um ângulo, sem precisar recalcular a soma $\theta_1 + \theta_2$. A partir desse ângulo, os dois dedos são desenhados com abertura de $\pm 0.6$ radianos, formando um "V" que acompanha o braço em qualquer configuração.
\end{itemize}

\section{Metodologia}
Nesta seção são descritos os materiais, a configuração do ambiente e a estrutura do código desenvolvido para a simulação.

\subsection{Materiais e Configuração}
O projeto foi desenvolvido em \textit{Python}, utilizando as bibliotecas \texttt{math}, \texttt{numpy}, \texttt{opencv-python} (\texttt{cv2}) e \texttt{matplotlib}, instaladas via:

\begin{lstlisting}[style=pythonstyle]
pip install opencv-python matplotlib numpy
\end{lstlisting}

A tela de renderização é alocada como uma matriz \texttt{numpy}, posteriormente manipulada pelo OpenCV:

\begin{lstlisting}[style=pythonstyle]
tela = np.zeros((400, 400, 3), dtype=np.uint8)
\end{lstlisting}

\subsection{Procedimentos e Scripts Utilizados}
O algoritmo principal foi modularizado em funções para separar a lógica de cálculo cinemático da lógica de renderização gráfica. As funções criadas são detalhadas a seguir.

A função \texttt{calcular\_junta} (Listagem \ref{lst:script_calculo}) implementa o cerne matemático, utilizando trigonometria básica (\texttt{math.cos} e \texttt{math.sin}) para definir as novas coordenadas, convertendo os resultados para inteiros, exigência das funções de desenho.

\begin{lstlisting}[style=pythonstyle, caption={Cálculo cinemático das juntas do braço robótico}, label={lst:script_calculo}]
def calcular_junta(x_origem, y_origem, comprimento, angulo_rad):
    x_final = x_origem + comprimento * math.cos(angulo_rad)
    y_final = y_origem - comprimento * math.sin(angulo_rad)
    return int(round(x_final)), int(round(y_final))
\end{lstlisting}

A função \texttt{desenhar\_engrenagem} (Listagem \ref{lst:script_engrenagem}) é uma melhoria estética. Ela desenha um círculo base e, iterativamente através da divisão dos 360 graus (ou $2\pi$), calcula e posiciona pequenos círculos (dentes) nas bordas do círculo principal, finalizando com um furo central escuro.

\begin{lstlisting}[style=pythonstyle, caption={Função estilizada para desenho de engrenagens utilizando aproximações trigonométricas}, label={lst:script_engrenagem}]
def desenhar_engrenagem(tela, centro_x, centro_y, raio, cor, num_dentes=8):
    """Desenha um formato de engrenagem estilizada"""
    # Desenha o circulo base da engrenagem
    cv2.circle(tela, (centro_x, centro_y), raio, cor, -1)
    
    # Desenha os dentes da engrenagem usando trigonometria
    for i in range(num_dentes):
        angulo = i * (2 * math.pi / num_dentes)
        
        # Calcula a ponta de cada dente
        dente_x = int(centro_x + (raio * 1.3) * math.cos(angulo))
        dente_y = int(centro_y + (raio * 1.3) * math.sin(angulo))
        
        # Desenha os dentes como pequenos circulos na borda
        cv2.circle(tela, (dente_x, dente_y), int(raio * 0.4), cor, -1)
        
    # Desenha um furo no meio (eixo preto/cinza escuro)
    cv2.circle(tela, (centro_x, centro_y), int(raio * 0.4), (50, 50, 50), -1)
\end{lstlisting}

A função \texttt{desenhar\_garra} (Listagem \ref{lst:script_garra}) é outra melhoria implementada. Ela calcula a orientação global do antebraço usando \texttt{math.atan2} para que a garra aponta sempre para a direção correta. A partir disso, dois dedos são gerados com uma abertura de $0.6$ radianos.

\begin{lstlisting}[style=pythonstyle, caption={Função estilizada para desenho da garra do robô}, label={lst:script_garra}]
def desenhar_garra(tela, cotovelo_x, cotovelo_y, mao_x, mao_y):
    """Desenha uma garra com 2 dedos apontando na direcao do antebraco"""
    # Calcula o angulo do antebraco para saber para onde a mao aponta
    angulo_braco = math.atan2(mao_y - cotovelo_y, mao_x - cotovelo_x)
    
    tamanho_dedo = 25
    abertura = 0.6 
    
    # Calcula as posicoes das pontas dos dedos esquerdo e direito
    ponta_esq_x = int(mao_x + tamanho_dedo * math.cos(angulo_braco - abertura))
    ponta_esq_y = int(mao_y + tamanho_dedo * math.sin(angulo_braco - abertura))
    
    ponta_dir_x = int(mao_x + tamanho_dedo * math.cos(angulo_braco + abertura))
    ponta_dir_y = int(mao_y + tamanho_dedo * math.sin(angulo_braco + abertura))

    # Desenha a base da mao e os dois dedos (linhas grossas)
    cor_garra = (255, 255, 0)
    cv2.circle(tela, (mao_x, mao_y), 10, cor_garra, -1) 
    cv2.line(tela, (mao_x, mao_y), (ponta_esq_x, ponta_esq_y), cor_garra, 5) 
    cv2.line(tela, (mao_x, mao_y), (ponta_dir_x, ponta_dir_y), cor_garra, 5) 
\end{lstlisting}

A função \texttt{desenhar\_braco\_estilizado} (Listagem \ref{lst:script_desenho_geral}) encarrega-se de sobrepor e unificar todas as formas geométricas (linhas, retângulos, engrenagens e garra) na matriz \texttt{tela}. Após compor a imagem, a biblioteca Matplotlib é utilizada para salvar a renderização final no arquivo \texttt{braco\_robotico\_estilizado.png}.

\begin{lstlisting}[style=pythonstyle, caption={Função unificadora para desenhar o braço completo e exportar a imagem}, label={lst:script_desenho_geral}]
def desenhar_braco_estilizado(base_x, base_y, cotovelo_x, cotovelo_y, mao_x, mao_y, tela):
    # Primitiva: Retangulo (a base fixa do robo)
    cv2.rectangle(tela, (base_x - 35, base_y), (base_x + 35, 400), (128, 128, 128), -1)

    # Primitiva: Linha (os elos rigidos)
    cv2.line(tela, (base_x, base_y), (cotovelo_x, cotovelo_y), (200, 200, 200), 8)
    cv2.line(tela, (cotovelo_x, cotovelo_y), (mao_x, mao_y), (200, 200, 200), 8)

    # Engrenagens
    desenhar_engrenagem(tela, base_x, base_y, raio=12, cor=(255, 0, 0))
    desenhar_engrenagem(tela, cotovelo_x, cotovelo_y, raio=10, cor=(255, 0, 0))

    # Garra
    desenhar_garra(tela, cotovelo_x, cotovelo_y, mao_x, mao_y)

    # Exibe a imagem
    plt.imshow(tela)
    plt.title("Braco Robotico Articulado (Estilizado)")
    plt.savefig("braco_robotico_estilizado.png")
\end{lstlisting}

A orquestração principal (Listagem \ref{lst:script_main}) gerencia a coleta dos ângulos escolhidos pelo usuário, aloca a tela virtual e chama o fluxo de renderização.

\begin{lstlisting}[style=pythonstyle, caption={Rotina principal de execução e coleta de dados}, label={lst:script_main}]
def definir_angulos():
    base_x, base_y = 120, 340
    comprimento_braco = float(input("Digite o comprimento do braco (em pixels) padrao 130: "))
    comprimento_antebraco = float(input("Digite o comprimento do antebraco (em pixels) padrao 100: "))
    ang_ombro = float(input("Digite o angulo do ombro (em graus) padrao 65: "))
    ang_cotovelo = float(input("Digite o angulo do cotovelo (em graus) padrao -40: "))

    ang_ombro_rad = math.radians(ang_ombro)
    ang_cotovelo_rad = math.radians(ang_cotovelo)
    return base_x, base_y, comprimento_braco, comprimento_antebraco, ang_ombro_rad, ang_cotovelo_rad

def main():
    print("Iniciando a projecao do braco robotico...")
    base_x, base_y, comprimento_braco, comprimento_antebraco, ang_ombro_rad, ang_cotovelo_rad = definir_angulos()
    tela = np.zeros((400, 400, 3), dtype=np.uint8)

    cotovelo_x, cotovelo_y = calcular_junta(base_x, base_y, comprimento_braco, ang_ombro_rad)
    mao_x, mao_y = calcular_junta(cotovelo_x, cotovelo_y, comprimento_antebraco, ang_ombro_rad + ang_cotovelo_rad)

    print(f"Ombro   --> Coordenadas: ({base_x}, {base_y})")
    print(f"Cotovelo --> Coordenadas: ({cotovelo_x}, {cotovelo_y})")
    print(f"Mao      --> Coordenadas: ({mao_x}, {mao_y})")

    desenhar_braco_estilizado(base_x, base_y, cotovelo_x, cotovelo_y, mao_x, mao_y, tela)

if __name__ == "__main__":
    main()
\end{lstlisting}

\section{Resultados e Análises}

Durante a execução da rotina \texttt{main()}, os parâmetros angulares e dimensionais podem ser definidos dinamicamente pelo usuário. Assumindo a entrada padrão:
\begin{itemize}
    \item Coordenadas da base: $(120, 340)$
    \item Comprimento do elo do ombro: $130$ px
    \item Comprimento do elo do antebraço: $100$ px
    \item Ângulo do ombro: $65^\circ$
    \item Ângulo relativo do cotovelo: $-40^\circ$
\end{itemize}

O sistema obteve êxito no cálculo progressivo de cada junta com base nas restrições trigonométricas do desafio. A matriz de transformação implícita propagou adequadamente as coordenadas cartesianas até a posição do efetuador final (garra).

As coordenadas projetadas no console, confirmam o alinhamento exato no plano computacional:
\begin{verbatim}
    Ombro   --> Coordenadas: (120, 340)
    Cotovelo --> Coordenadas: (175, 222)
    Mão      --> Coordenadas: (266, 180)
\end{verbatim}

A renderização visual resultante pode ser conferida na Figura \ref{fig:braco_estilizado}, ilustrando o mapeamento de coordenadas matemáticas para os índices da matriz de pixels, respeitando o espaço de tela com a origem no quadrante superior esquerdo.

\begin{figure}[H]
    \centering
    \includegraphics[width=0.9\textwidth]{images/braco_robotico_estilizado.png}
    \caption{Renderização final do braço robótico estilizado, apresentando a base, os braços e as melhorias aplicadas às juntas (engrenagens) e extremidade (garra).}
    \label{fig:braco_estilizado}
\end{figure}

\section{Conclusão}

O experimento atingiu amplamente os objetivos do Desafio do Braço Robótico. Foi realizada a conversão bem-sucedida de conceitos matemáticos abstratos de cinemática e trigonometria em uma representação gráfica bidimensional estruturada em \textit{Screen Space}. As adições de componentes estilizados demonstraram que o domínio matricial permite o desenho de mecanismos complexos partindo de primitivas básicas.

\begin{thebibliography}{9}

\bibitem{rodrigo_primitivas}
RODRIGO, Raulivan. \textbf{Desenhar primitivas: Computação Gráfica}. Notas de Aula (05 - Desenhar primitivas.pdf). Departamento de Computação, Centro Federal de Educação Tecnológica de Minas Gerais (CEFET-MG), 2026.

\end{thebibliography}

\end{document}
